\section*{الفصل \ref{ch:g12:past-climates} --- قراءة المناخات الماضية}

\begin{solution}{exo:g12:past-climates:1}
(أ) \omterm{def:g12:past-climates:proxy}{مؤشر بديل}؛ (ب) \omterm{def:g12:past-climates:ice-core}{عينة جليدية}؛ (ج) \omterm{def:g12:past-climates:glacial-cycle}{فترة بين جليدية}؛ (د) \omterm{def:g12:past-climates:pollen}{تحليل حبوب اللقاح}؛
(ه) \omterm{def:g12:past-climates:lgm}{الذروة الجليدية الأخيرة}.
\end{solution}

\begin{solution}{exo:g12:past-climates:2}
(أ) الجليد القطبي: ثاني أكسيد الكربون والميثان في هواء ذلك الزمن.
(ب) رسوبيات البحيرات أو الخث: الغطاء النباتي، وهنا أشجار تحتاج إلى صيف
دافئ. (ج) رسوبيات المحيطات: حجم الجليد على اليابسة ودرجة حرارة
\omterm{def:g10:layered-ocean:thermocline}{المحيط العميق}. (د) الأشجار: ظروف النمو، سنة بسنة.
\end{solution}

\begin{solution}{exo:g12:past-climates:3}
(أ) $-50 - (-120) = \qty{70}{m}$ في \num{7000} سنة: $\qty{70000}{mm}/\num{7000}
= \qty{10}{mm}$ في السنة. (ب) $10/3.7 \approx 2.7$: كان متوسط ارتفاع البحر في
انحسار الجليد أسرع بنحو 2.7 مرة من ارتفاعه في 2006--2018.
\end{solution}

\begin{solution}{exo:g12:past-climates:4}
القرن عشرة عقود: $\qty{2.0}{\degreeCelsius}$ لكل قرن. والألف سنة
مئة عقد: $\qty{20}{\degreeCelsius}$ لكل ألف سنة.
\end{solution}

\begin{solution}{exo:g12:past-climates:5}
(أ) \num{23000} سنة: مبادرة المحور؛ و\num{42000} سنة: ميل
المحور؛ ونحو \num{100000} سنة: شكل المدار.
(ب) $\num{1000000}/\num{41000} \approx 24$ دورة؛ و$\num{1000000}/\num{100000}
= 10$ دورات.
\end{solution}

\begin{solution}{exo:g12:past-climates:6}
(أ) $266.6 + 7.83 \times (-9.4) \approx 266.6 - 73.6 = \qty{193}{ppm}$: أي
أعلى بمقدار \qty{3}{ppm} فقط من القيمة المقيسة \qty{190}{ppm}. (ب) $266.6 + 7.83 \times 5.5
\approx \qty{310}{ppm}$، أي أعلى بمقدار \qty{11}{ppm} من أعلى قيمة مقيسة: فأدفأ
نقطة تقع تحت الخط. (ج) $\Delta T = (427 - 266.6)/7.83 \approx
\qty{20}{\degreeCelsius}$. فالخط لا يصف إلا كيف تغيرت الكميتان
معًا، في دورات طبيعية، داخل مدى البيانات
(من \qty{-11}{\degreeCelsius} إلى \qty{+6}{\degreeCelsius}، ومن 174 إلى
\qty{299}{ppm})؛ وهو ليس قانونًا، ولا يمكن أن تُقرأ منه قيمة بعيدة
جدًا عن ذلك المدى.
\end{solution}

\begin{solution}{exo:g12:past-climates:7}
(أ) الأعشاب، \qty{84}{\percent}: منظر مفتوح يكاد يخلو من الأشجار، بارد.
(ب) البلوط (\qty{42}{\percent}) والبندق (\qty{34}{\percent}): غابة عريضة
الأوراق، صيفها دافئ. (ج) ينتشر البتولا أولًا، ثم الصنوبر؛ ويصل البندق والبلوط
أخيرًا. (د) لا تنتشر الأشجار إلا بسرعة انتقال بذورها ونمو
صغارها، انطلاقًا من الأماكن التي نجت فيها من البرد؛ ويجب أيضًا أن تتكوّن
تربة. فقد تتأخر الغابة عن المناخ.
\end{solution}

\begin{solution}{exo:g12:past-climates:8}
(أ) \qty{300}{cm}، أي \qty{3000}{mm}، في \num{8000} سنة: $3000/8000 =
\qty{0.375}{mm}$ في السنة. (ب) $\num{6000} + (470 - 300)/300 \times \num{8000}
\approx \num{6000} + \num{4530} \approx \num{10500}$ سنة. (ج) $10/0.375
\approx 27$ سنة.
\end{solution}

\begin{solution}{exo:g12:past-climates:9}
(أ) $7.16 - 4.11 = 3.05$ مليون كيلومتر مربع، و$3.05/4.11 \approx
\qty{74}{\percent}$ من اليابسة الإضافية. (ب) نعم: ففي \omterm{def:g12:past-climates:lgm}{الذروة الجليدية الأخيرة} كان البحر
أدنى من مستواه اليوم بمقدار 125 إلى \qty{134}{m}، وهذا أكثر بكثير من \qty{36}{m}
التي تكفي لوصل بورنيو بآسيا على الخريطة. (ج) تستعمل الخريطة قاع البحر الحالي، الذي لم
يبق على حاله تمامًا منذ ذلك الحين (فقد ترسبت عليه رسوبيات، وتحركت
الأرض قليلًا)؛ والشبكة أيضًا خشنة، نقطة واحدة كل 0.1
درجة، ولا يُعرف انخفاض البحر إلا في حدود نحو عشرة أمتار.
\end{solution}

\begin{solution}{exo:g12:past-climates:10}
(أ) قيمة \omterm{def:g12:past-climates:lgm}{الذروة الجليدية الأخيرة}، \qty{4.92}{\permil}، هي
الأعلى. وهذا متوقع: فالأغطية الجليدية الكبيرة في ذلك الزمن خزنت ماءً فقيرًا
بالأكسجين-18، فغادر المحيط، وترك المحيط والأصداف التي نمت فيه أغنى
بالأكسجين-18؛ وكان الماء العميق أبرد أيضًا. (ب) كي تُرسم الأزمنة الباردة الجليدية
في الأسفل والأزمنة الدافئة في الأعلى، كما في منحنى درجة الحرارة.
\end{solution}

\begin{solution}{exo:g12:past-climates:11}
(أ) يُظهر الارتباط أن السجلين يرتفعان وينخفضان معًا، قريبًا
من خط مستقيم؛ ولا يُظهر أيهما سبّب الآخر، ولا أن أحدهما
سبّبه أصلًا. (ب) ربما أطلق مناخ أدفأ ثاني أكسيد الكربون في الهواء؛
أو ربما حرّك سبب ثالث، كتغيرات مدار الأرض،
الاثنين معًا؛ أو ربما دفع كل منهما الآخر. (ج) توقيت التغيرات، أي
أيها جاء أولًا، بتأريخ دقيق للجليد ولهوائه؛ وفهم
كيف يؤثر ثاني أكسيد الكربون في إشعاع الأرض
(\cref{ch:g12:climate-change}).
\end{solution}

\begin{solution}{exo:g12:past-climates:12}
(أ) $\num{1300000}/\num{41000} \approx 32$ دورة بين 2.6 و1.3
مليون سنة مضت؛ و$\num{700000}/\num{98500} \approx 7$ دورات على مدى
\num{700000} السنة الأخيرة. (ب) \num{41000} سنة: ميل المحور؛ ونحو
\num{100000} سنة: شكل المدار. (ج) الدورة البالغة \num{41000} سنة
تُقاس عيناتها نحو $\num{41000}/\num{2000} \approx 20$ مرة، وهذا يكفي لرسمها؛
أما التغير الذي يدوم قرنًا فيقع بين عينتين، أو داخل عينة واحدة،
ولا يكاد يترك أثرًا.
\end{solution}

\begin{solution}{exo:g12:past-climates:13}
(أ) $192/10 \approx 19$ و$198/10 \approx 20$: أسرع بنحو 20 مرة.
(ب) مقارنةً بمعدل \qty{4}{ppm} لكل قرن: من $192/4 = 48$ إلى $198/4 \approx 50$
مرة؛ ومقارنةً بمعدل \qty{42}{ppm}: من $192/42 \approx 4.6$ إلى $198/42 \approx 4.7$
مرة. فالارتفاع اليوم أسرع بنحو 5 إلى 50 مرة من ارتفاع الحد الأقصى الحراري.
(ج) الارتفاع $2000 - 900 = \qty{1100}{ppm}$: أي $1100/42 \approx 26$ قرنًا،
نحو \num{2600} سنة بأسرع معدلات الحد الأقصى الحراري؛ و$1100/195 \approx 5.6$
قرن، أي أقل من 600 سنة، بمعدل اليوم.
\end{solution}

\begin{solution}{exo:g12:past-climates:14}
(أ) عشرة أضعاف 1 إلى \qty{1.5}{\degreeCelsius}: من 10 إلى \qty{15}{\degreeCelsius}
لكل ألف سنة، أي من 0.10 إلى \qty{0.15}{\degreeCelsius} لكل عقد.
(ب) المعدل العالمي اليوم، \qty{0.20}{\degreeCelsius} لكل عقد، يساوي من 1.3 إلى 2
ضعف أسرع احترار فوق غرينلاند. (ج) يمكن لمكان واحد أن يتغير أكثر بكثير
من المتوسط على الكرة الأرضية كلها، حيث تتعادل تغيرات المناطق المختلفة
جزئيًا: فالمعدل العالمي يُقارن بمعدل عالمي.
\end{solution}

\begin{solution}{exo:g12:past-climates:15}
(أ) \qty{14}{m} بمعدل \qty{12}{m} لكل ألف سنة تستغرق $14/12 \approx
1.17$ ألف سنة، أي نحو \num{1170} سنة: أي قبل نحو \num{9800} سنة.
(ب) لا: فقد كان على البحر أن يبقى أدنى من مستواه اليوم بمقدار \qty{36}{m} على الأقل ليصل
بورنيو بآسيا؛ وحين تجاوز \qty{-36}{m} غمر الماء الجسر البري،
وعادت بورنيو جزيرة. (ج) $36/12 = 3$ آلاف سنة، أي نحو
\num{3000} سنة؛ والواقع أن الارتفاع تباطأ بعد نحو \num{8200} سنة
مضت، فاستغرق مدة أطول.
\end{solution}

\begin{solution}{pb:g12:past-climates:1}
\textbf{1.} الجليد نفسه، بنظائره الثقيلة، يسجل
درجة الحرارة؛ والهواء في فقاعاته يسجل ثاني أكسيد الكربون.

\textbf{2.} $5.5 - (-10.6) = \qty{16.1}{\degreeCelsius}$.

\textbf{3.} $299 - 174 = \qty{125}{ppm}$.

\textbf{4.} تقع النقاط قريبًا من خط مستقيم: فقد ارتفع السجلان
وانخفضا معًا. ولا يقول أيهما سبّب الآخر.

\textbf{5.} تُرسم سجلات الماضي والأزمنة الأقدم على أحد
الجانبين أو الآخر؛ فإذا قُرئ السجل في الاتجاه الخطأ، حُسب الاحترار تبردًا.

\textbf{6.} \omterm{def:g12:past-climates:glacial-cycle}{الفترة الجليدية} فترة باردة داخل عصر جليدي، تنمو فيها
الأغطية الجليدية إلى ما يتجاوز حجمها الحالي بكثير؛ \omterm{def:g12:past-climates:glacial-cycle}{والفترة بين الجليدية} فترة دافئة
بين \omterm{def:g12:past-climates:glacial-cycle}{فترتين جليديتين}، تكون فيها الأغطية الجليدية قريبة من حجمها اليوم أو
أصغر.

\textbf{7.} في \omterm{def:g12:past-climates:glacial-cycle}{فترة بين جليدية}، بدأت قبل نحو \num{11000} سنة.

\textbf{8.} $\num{800000}/\num{98000} \approx 8$: نحو ثماني دورات.

\textbf{9.} شكل المدار.

\textbf{10.} إيقاع دورته \num{41000} سنة، وهي دورة ميل
المحور.

\textbf{11.} $-0.4 - (-9.4) = \qty{9.0}{\degreeCelsius}$.

\textbf{12.} $264 - 190 = \qty{74}{ppm}$.

\textbf{13.} \qty{6}{\degreeCelsius}؛ وتغير القارة القطبية الجنوبية، \qty{9.0}{\degreeCelsius}،
أكبر، نحو 1.5 ضعف التغير العالمي.

\textbf{14.} $\num{18000} - \num{11000} = \num{7000}$ سنة.

\textbf{15.} $6/7 \approx \qty{0.86}{\degreeCelsius}$ لكل ألف سنة.

\textbf{16.} الألف سنة مئة عقد: $0.86/100 \approx
\qty{0.0086}{\degreeCelsius}$ لكل عقد.

\textbf{17.} $0.20 \times 100 = \qty{20}{\degreeCelsius}$ لكل ألف سنة.

\textbf{18.} $20/0.86 \approx 23$ مرة.

\textbf{19.} $20/1.5 \approx 13$ و$20/1 = 20$: من 13 إلى 20 مرة.

\textbf{20.} الاحترار العالمي اليوم \textbf{أسرع بنحو عشرين مرة} من
الاحترار الخارج من آخر عصر جليدي: 23 مرة من متوسط معدله، ومن 13 إلى 20 مرة
من أسرع معدلاته.
\end{solution}
